%% REWRITTEN 2026-10-16. The old version showed the k-slot rotation but not the
%% thing that makes it necessary: the per-layer deadline, and what happens to the
%% slack as k changes. A reader could see the mechanism and not the reason.
\begin{figure}[t]
\centering
\resizebox{\columnwidth}{!}{%
\begin{tikzpicture}[
  x=1cm, y=1cm, font=\scriptsize,
  cmp/.style={draw, fill=black!8, minimum height=5.5mm, inner sep=1pt},
  xfer/.style={draw, fill=blue!15, draw=blue!50, minimum height=4.5mm, inner sep=1pt},
  late/.style={draw, fill=red!18, draw=red!55, minimum height=4.5mm, inner sep=1pt},
  lbl/.style={font=\scriptsize\itshape, text=black!60, anchor=east},
  note/.style={font=\tiny, text=black!65, align=left},
  hd/.style={font=\bfseries\scriptsize, anchor=west},
  dl/.style={red!60, densely dashed},
]

%% ================= (a) the deadline, and k=1 ==========================
\node[hd] at (-1.9,2.5) {(a) the deadline: layer $\ell$ cannot start until its bytes have landed};

\node[lbl] at (-0.15,1.55) {compute};
\node[lbl] at (-0.15,0.80) {link};

\node[cmp, minimum width=1.5cm, anchor=west] at (0,1.55) {layer $\ell$};
\node[late, minimum width=1.9cm, anchor=west] at (1.55,0.80) {bytes for $\ell{+}1$};
\node[cmp, minimum width=1.5cm, anchor=west] at (3.50,1.55) {layer $\ell{+}1$};

\draw[dl] (1.55,1.95) -- (1.55,0.45);
\draw[dl] (3.45,1.95) -- (3.45,0.45);
\node[note, anchor=south, text=red!70] at (2.50,1.98) {transfer is \emph{exposed}};
\draw[{Latex[length=1.4mm]}-{Latex[length=1.4mm]}, red!60] (1.55,0.35) -- (3.45,0.35);
\node[note, anchor=north, text=red!70, align=center] at (2.50,0.28)
     {the step simply waits};

\node[note, anchor=west] at (5.8,1.2)
     {With one slot there is nowhere to put layer $\ell{+}1$'s bytes\\
      until layer $\ell$ has finished reading them. The transfer\\
      therefore cannot begin early, and its whole duration is\\
      added to the step.};

%% ================= (b) k slots give slack =============================
\begin{scope}[yshift=-3.3cm]
\node[hd] at (-1.9,2.5) {(b) $k$ slots: a slot is not reused for $k$ layers, so a transfer may start $k{-}1$ layers early};

\node[lbl] at (-0.15,1.55) {compute};
\node[lbl] at (-0.15,0.80) {link};

\foreach \i/\x/\t in {0/0/{$\ell$}, 1/1.55/{$\ell{+}1$}, 2/3.10/{$\ell{+}2$},
                      3/4.65/{$\ell{+}3$}, 4/6.20/{$\ell{+}4$}}
  \node[cmp, minimum width=1.5cm, anchor=west] at (\x,1.55) {layer \t};

\foreach \x/\t in {0/{$\ell{+}4$}, 1.55/{$\ell{+}5$}, 3.10/{$\ell{+}6$},
                   4.65/{$\ell{+}7$}, 6.20/{$\ell{+}8$}}
  \node[xfer, minimum width=1.5cm, anchor=west] at (\x,0.80) {\t};

\node[note, anchor=west] at (8.0,1.2)
     {Each transfer now has $k{-}1$ layers of compute to\\
      finish under. It is hidden entirely whenever\\
      \ \ \ \ \textbf{bytes} / \textbf{link rate} $\;\le\;$ one layer's compute,\\
      and that inequality --- not $k$ --- is what decides\\
      how many host-tier requests a step can carry.};

%% slot reuse bracket
\draw[decorate, decoration={brace, amplitude=1.2mm}, black!45]
     (0,0.45) -- (6.20,0.45);
\node[note, anchor=north, text=black!60] at (3.10,0.30)
     {slot $s_j$ is written here and not read until $k$ layers later};
\end{scope}

%% ================= (c) what k does and does not buy ===================
\begin{scope}[yshift=-5.15cm]
\node[hd] at (-1.9,1.35) {(c) what more slots do, and do not, buy};
\node[draw, rounded corners=2pt, fill=black!3, align=left, inner sep=2mm,
      anchor=north west, font=\scriptsize] at (-1.9,1.05)
     {\textbf{more slack} against jitter: a late transfer has more layers to catch up in.\\[2pt]
      \textbf{not more bandwidth}: the hiding condition above contains no $k$. Going from\\
      $k{=}2$ to $k{=}8$ moves neither coefficient measurably, while every extra slot\\
      costs device memory --- which is the resource the whole design exists to save.};
\end{scope}

\end{tikzpicture}}
\caption{\textbf{Why the transfer is issued ahead, and what the slot count is for.}
Each layer's attention cannot begin until that layer's host-resident KV is on the
device, so with a single staging slot the transfer is serialised against compute and
fully exposed (a). Holding $k$ slots means a slot is not overwritten for $k$ layers,
so a transfer can be issued that far ahead and finish underneath intervening compute
(b). The slot count buys \emph{tolerance}, not throughput: whether a transfer hides
at all depends on the per-layer byte time against the per-layer compute time, an
inequality in which $k$ does not appear (c).}
\label{fig:pipeline}
\end{figure}
